\section{Block Composition, Strong Powers, and Asymptotic Capacity}\label{sec:capacity-theory}
%==============================================================================


This section studies how the induced failure topology scales. Once exact recovery is governed by colorability of the one-shot confusability graph, repeated composition does not reset the ambiguity; it composes the same failure structure across blocks. The next results identify the correct strong-power object and the resulting asymptotic zero-error rate.

\subsection{Motivation: Repeated Composition}

What happens if we repeat the partial-view experiment $n$ times? If we compose $n$ independent copies of the same encoding system, does the ambiguity explode uncontrollably, or does it grow at a predictable rate?

The answer depends on the failure topology:
\begin{itemize}
\item If the system were \emph{fully} incoherent (confusability graph is a clique), every block would multiply the confusion. The ambiguity would grow without any exploitable structure.
\item Because our confusability graph has \emph{structure} (as in the 4-cycle example of Section~\ref{sec:binary-square}), the ambiguity growth is controlled by that structure. Some state pairs remain distinguishable even across multiple blocks.
\end{itemize}

Under $n$-fold block composition, the confusability graph becomes its $n$-fold \emph{strong power}:
\[
G_n \;=\; G^{\boxtimes n}.
\]
Independent sets in the strong power relate systematically to independent sets in the base graph, yielding a supermultiplicative growth law that converges to a well-defined asymptotic rate. The theorem below quantifies this rate.

\subsection{Block Composition Law}

A genuine supermultiplicative block law holds:
\[
\alpha_{m+n}\;\ge\;\alpha_m\alpha_n,
\]
where $\alpha_n$ denotes the largest one-shot zero-error code size in the $n$-block composed system. This is the correct discrete precursor to Shannon-capacity analysis: the admissible code-size sequence is already supermultiplicative at the graph level.\leanmeta{\LH{MFT22}}

More sharply, the block-level graph itself factors correctly. Confusability between concatenated $(m+n)$-block states is equivalent to adjacency in the strong product of the $m$-block and $n$-block confusability graphs, packaged formally as an explicit graph isomorphism. Thus the repeated-block system is not only bounded by strong-product arguments at the codebook level; it is literally a strong-product graph-power object for the same induced failure map.\leanmeta{\LHrng{MFT}{29}{30}, \LHrng{MFT}{38}{39}}

\paragraph{Conjunction, not disjunction.}
Block confusability requires a \emph{single location tuple} $\ell = (\ell_1, \ell_2)$ that works \emph{simultaneously} for both blocks:
\begin{itemize}
\item Block 1: observe location $\ell_1$, must match
\item Block 2: observe location $\ell_2$, must match
\end{itemize}
Two $(m+n)$-block states are confusable iff there exists one location tuple making observations match across \emph{all} blocks. This is not ``confusable in block 1 OR confusable in block 2'' but rather ``$\exists$ location tuple, $\forall$ blocks: match.'' That universal quantification over blocks is exactly the strong product condition: for each block, either equal or confusable (matching on some location), with at least one block strictly confusable.

\paragraph{Concrete example.}
Consider two copies of the binary square. State $((0,0), (0,0))$ vs $((0,1), (1,1))$:
\begin{itemize}
\item Block 1: $(0,0) \sim (0,1)$ in $C_4$ (agree on $x_1$)
\item Block 2: $(0,0) \not\sim (1,1)$ in $C_4$ (disagree on both coordinates)
\end{itemize}
An ``OR-product'' would mark these confusable because block 1 matches. But the actual model requires \emph{one} location pair $(\ell_1, \ell_2)$ matching \emph{both} blocks. Block 1 can match via $\ell_1 = 1$ (observing $x_1$). Block 2 cannot match: location 1 gives $0$ vs $1$, location 2 gives $0$ vs $1$. No location pair works, so these states are not confusable, exactly as the strong product predicts.

Iterating the same construction yields the derived $n$-block lower bound. If $\alpha_1$ denotes the one-shot maximum independent-set code size, then the mechanized block-power theorem gives
\[
\alpha_1^n \;\le\; \alpha_n,
\]
where $\alpha_n$ is the largest one-shot zero-error code size in the $n$-fold block-composed system. This is the first asymptotic-capacity scaffold in the extended model: repeated block composition already forces the expected exponential growth law before any Shannon-capacity or $\vartheta$-style upper theory is introduced. Once the architecture determines the one-shot failure graph, that same graph already controls the asymptotic lower law.\leanmeta{\LH{MFT21}}

For notation, write $\Theta(G)$ for the Shannon capacity of a confusability graph $G$ in logarithmic units, and write $\vartheta_\infty(G)$ for the corresponding asymptotic Lov\'asz-$\vartheta$ upper value. When the object is a view family $\mathcal V$ rather than an abstract graph, we write $\Theta(G_{\mathcal V})$ and $\vartheta_\infty(G_{\mathcal V})$ for the same quantities computed on the induced confusability graph $G_{\mathcal V}$.

These normalized block rates can be packaged into an explicit lower asymptotic capacity envelope
\[
C_* \;:=\; \sup_{n\ge 1} \frac{\log \alpha_n}{n},
\]
implemented as an $\mathrm{ENNReal}$ supremum of the positive block-rate sequence. Formally, $\alpha_n$ is identified directly with the maximum finite independent-set size of the $n$-block confusability graph, the block graph is identified with the $n$-fold strong power of the one-shot confusability graph, and the same envelope is proved equal to the graph-side Shannon lower-capacity expression built from those strong powers.\leanmeta{\LHrng{MFT}{31}{32}, \LHrng{MFT}{35}{36}}

The asymptotic scaffold is also not limited to a lower envelope statement. On the graph side itself, strong-product lower bounds for independent-set cardinality induce monotonicity of normalized strong-power rates along repeated multiples. Specializing those generic graph theorems back to the model-generated confusability graph yields the repeated-block statement below.\leanmeta{\LHrng{MFT}{40}{42}}

Repeating a fixed $m$-block system $k$ times yields
\[
\alpha_m^k \;\le\; \alpha_{km},
\]
and hence
\[
\frac{\log \alpha_m}{m} \;\le\; \frac{\log \alpha_{km}}{km}.
\]
Thus rates along repeated block multiples are monotone from below toward the same capacity envelope. This is still a lower theory rather than a full Shannon-capacity theorem, but it is already a genuine asymptotic graph-rate statement for the induced failure topology.\leanmeta{\LH{MFT33}}

Once the block confusability graph has been identified with strong powers of the one-shot graph, the asymptotic rate theory is exactly the classical Shannon-capacity framework specialized back to the model-generated graph family.

\begin{theorem}[Asymptotic Shannon-capacity theorem]\label{thm:asymptotic-capacity}
Let $\alpha_n$ denote the largest one-shot zero-error code size in the $n$-block composed system. Then the normalized block rates
\[
\frac{\log \alpha_n}{n}
\]
converge to a real asymptotic Shannon-capacity value, and that value is equal to the supremum of the finite block-rate sequence.
\end{theorem}
\leanmeta{\LHrng{MFT}{43}{49}, \LH{MFT56}}

\begin{proof}
The argument rests on two observations: supermultiplicativity of independent-set sizes under strong product, and Fekete's lemma.

\emph{Supermultiplicativity.} Recall that the $n$-block confusability graph $G_n$ is the $n$-fold strong power $G^{\boxtimes n}$ of the one-shot graph $G$. For any two graphs $H$ and $K$, an independent set in $H\boxtimes K$ can be constructed from independent sets in $H$ and $K$ as follows. If $I_H\subseteq V(H)$ and $I_K\subseteq V(K)$ are independent, then the Cartesian product $I_H\times I_K\subseteq V(H)\times V(K)$ is independent in $H\boxtimes K$: if $(u,v)$ and $(u',v')$ are adjacent in the strong product, then either $u=u'$ and $v\sim v'$, or $u\sim u'$ and $v=v'$, or $u\sim u'$ and $v\sim v'$. In each case at least one of $u,u'$ or $v,v'$ is an edge within the respective independent set, which is impossible. Hence $\alpha(H\boxtimes K)\ge\alpha(H)\alpha(K)$. Iterating gives $\alpha_{m+n}=\alpha(G^{\boxtimes(m+n)})\ge\alpha(G^{\boxtimes m})\alpha(G^{\boxtimes n})=\alpha_m\alpha_n$, so the sequence $\{\alpha_n\}$ is supermultiplicative.

\emph{Fekete convergence.} The inequality $\alpha_{m+n}\ge\alpha_m\alpha_n$ is equivalent to subadditivity of the sequence $\{-\log\alpha_n\}$. The trivial bound $\alpha_n\le|V(G)|^n$ ensures that $-\log\alpha_n\ge -n\log|V(G)|$, so the sequence is bounded below. By Fekete's subadditivity lemma, the normalized limit
\[
\lim_{n\to\infty}\frac{\log\alpha_n}{n}=\sup_{n\ge1}\frac{\log\alpha_n}{n}
\]
exists and equals the supremum of the finite block-rate sequence. This limit is the asymptotic Shannon capacity $\Theta(G)$ in logarithmic units.

The mechanized development packages the same argument with explicit sequence objects and convergence lemmas; the textual expansion above makes the independent-set product construction and Fekete application explicit for accessibility.
\end{proof}

\paragraph{Running example: Binary square.}
For the 4-cycle $C_4$ from Section~\ref{sec:binary-square}, the independence number is $\alpha(C_4) = 2$ (choose either diagonal). The strong power $C_4^{\boxtimes n}$ has $\alpha(C_4^{\boxtimes n}) = 2^n$, giving normalized rate $(\log 2^n)/n = \log 2$. The asymptotic Shannon capacity is therefore $\Theta(C_4) = \log 2$. This matches the exact tag-budget growth from the iterated binary-square family discussed earlier: the one-shot budget is $2$, and the $m$-block exact budget grows as $2^m$.

\paragraph{Relation to graph capacity theory.}
This places the extended model directly in the classical zero-error graph-capacity setting. The difference from the classical channel picture is generative rather than formal: here the confusability graphs arise from deterministic view families and partial observations on latent tuples. We do not claim a new general Shannon-capacity theorem for arbitrary graphs; the new content is that this partial-view model reaches the classical graph machinery in a non-clique regime, so the architecture determines the graph once and that graph then determines both one-shot and asymptotic recoverability. The clique-fiber regime recovers the classical cluster-graph equality case, while the transitivity criterion identifies when that case is produced by the view-family model. The non-clique witness already shows that the model also generates genuinely nontrivial graph behavior.
